\section{Introduction}
\label{sec:intro}

Variational Monte Carlo (VMC) approximates a quantum ground state by optimizing a parametrized trial wavefunction against expectation values estimated from sampled configurations. Restricted Boltzmann machines (RBMs) are an attractive choice of trial state because their amplitudes and logarithmic derivatives factorize over a hidden layer and can be evaluated in closed form~\cite{Carleo2017,Sorella2001}. Their practical accuracy, however, depends on three largely independent things: whether the sampler used during training actually draws from the state's Born distribution, whether the network's connectivity and sign structure can represent the target state at all, and whether the reported convergence is confirmed by an evaluation that does not share the training sampler's own biases.

A quantum annealer can supply the configurations used in this loop: a programmable Ising energy is realized on the device, and repeated anneals return classical spin strings~\cite{Gardas2018}. Whether those strings behave like Gibbs samples at some effective temperature is an empirical question, not a guarantee of the hardware. Benedetti \textit{et al.} showed that this effective temperature is instance dependent~\cite{Benedetti2016}; Nelson \textit{et al.} examined the conditions under which annealer output resembles thermal sampling at all~\cite{Nelson2022}. Any procedure that uses annealer samples inside an RBM training loop therefore needs an online, per-instance estimate of that effective temperature, obtained without extra hardware calls and without assuming the partition function is tractable.

Two recent papers make this a crowded claim on their own terms, and the present contribution must be stated more narrowly than ``online calibration from training samples'' as a result. Goto and Ohzeki already reuse visible and hidden samples generated during RBM training to calibrate more than one internal annealer parameter online, benchmarking directly against Gibbs sampling and contrastive divergence~\cite{Goto2025}. Kim, Gyhm, and Park go further and report that a diabatic-annealing sampler, corrected by an analytical temperature-rescaling method, gives RBM training with faster convergence and lower validation error than classical sampling outright~\cite{Kim2026}. What we isolate below is not ``online, sample-reusing calibration'' in general -- both of the above already do that -- but the specific \emph{pooled, no-requery} construction of Sec.~\ref{sec:cem-ls} and a deliberately narrower empirical claim: Sec.~\ref{sec:limitations} reports one criterion-specific, statistically marginal ordering rather than the general classical-vs-quantum superiority Kim \textit{et al.} assert.

\subsection{Pooled conditional thermometry: the central contribution}

Kubo and Goto introduced conditional expectation matching (CEM): the hidden-unit conditional means of an RBM, $\mathbb E[h_j\mid v]=\tanh(\beta\phi_j(v))$, are fit to the corresponding sampled hidden bits at a \emph{fixed} visible configuration, giving a partition-function-free estimate of $\beta$~\cite{Kubo2025,Kubo2026}. The contribution we isolate and develop here is the \emph{pooled} version of this idea: rather than querying the sampler again at a fixed $v$, every visible configuration $v^{(r)}$ produced during ordinary training is treated as its own conditioning vector, and all $(v^{(r)},h^{(r)})$ pairs already generated by stochastic reconfiguration (SR) are pooled into one least-squares (Sec.~\ref{sec:cem-ls}) estimate of a single effective inverse temperature $\betaeff$. This requires no additional annealer queries and no evaluation of $Z_\theta$.

We are explicit about what this estimator does and does not certify. A perfect fit of the hidden conditional means at a given $\beta$ is compatible with an arbitrarily wrong visible marginal (Sec.~\ref{sec:cem-limits}); consequently, a pooled-CEM estimate is a diagnostic of the hidden-layer response scale, not a certificate that the sampler reproduces $\pi_\theta(v)=|\psi_\theta(v)|^2$. Separating these two claims -- and building an independent-evaluation pipeline that tests the second one directly -- is the paper's second contribution.

\subsection{Related constraints: representation and validation}

Two further limitations are independent of thermometry and are treated as such rather than folded into a single ``competitiveness'' narrative. First, connectivity: a dense RBM does not fit natively on an annealer's coupler graph, so either minor embedding or an explicitly sparse, hardware-native mask must be used, and the two have different accuracy-vs-resource trade-offs~\cite{LeRoux2008}. Second, sign structure: a positive-amplitude RBM can only represent target states whose off-diagonal Hamiltonian matrix elements are non-positive in the computational basis~\cite{Bravyi2008}; the Marshall transformation supplies this sign for a bipartite nearest-neighbor antiferromagnet~\cite{Marshall1955} but not for a frustrated one~\cite{Richter1994,Voigt2000}. Neither limitation is repaired by a better sampler, and neither is the subject of this paper's central claim; we report both only insofar as they interact with the resource accounting in Sec.~\ref{sec:resources}.

\subsection{Enumerated contributions}

\begin{enumerate}
\item A pooled conditional-thermometry estimator for annealer-sampled RBMs, built on the least-squares objective of prior work (Sec.~\ref{sec:thermometry}), together with an explicit statement of what agreement of hidden conditional means does and does not imply about the visible distribution (Sec.~\ref{sec:cem-limits}).
\item An independent validation protocol -- exact enumeration at small $N$, long independent reference sampling at larger $N$, with training and evaluation using disjoint random streams -- applied to frozen, checkpointed ansätze, so that a reported convergence is confirmed by an estimator that does not share the training sampler's bias (Sec.~\ref{sec:validation}).
\item A quantitative account of the residual gap between the fitted effective temperature and the true visible distribution, on transverse-field-Ising systems small enough for exact enumeration (Sec.~\ref{sec:calibration}).
\item Independently validated time-to-accuracy curves across a grid of convergence tolerances $\epsilon$ and fields $h$, replacing a single-threshold, single-field time-to-$\epsilon$ statistic with its criterion dependence made explicit (Sec.~\ref{sec:tta}).
\item A joint accounting of device time, end-to-end wall time, and quantum-annealer hardware resources (physical qubits, chain length, chain-break fraction, embedding success rate) used to reach that accuracy, so that a timing comparison is read together with the resources it consumed (Sec.~\ref{sec:resources}).
\item An explicit statement of the conditions under which a quantum-annealer/FPGA crossover is observed, replacing an unconditional claim with the size-, field-, and tolerance-dependent statement the data actually support (Sec.~\ref{sec:limitations}).
\end{enumerate}

This paper is deliberately narrower than the research-stay report it is drawn from. Sparse hardware-native RBMs, parallel embeddings, and the Marshall-sign study are retained only as short, explicitly scoped appendices (App.~\ref{app:sparse}--\ref{app:marshall}); they support the resource-accounting discussion but are not part of the central thermometry-and-validation claim.
